%% sections/e-ramseyer.tex - H. R. 9510 Bill v4.0 (full bill) - Appendix E
%% Genre: Ramseyer comparative print (redline). Source deliverable 05.
%% Existing text plain; deletions \cprstrike{}; insertions \cprinsert{}.
\billsec{APPENDIX E.}{RAMSEYER COMPARATIVE PRINT (NON-OPERATIVE).}\phantomsection\label{toc:aE}

\uscprov{1}{\textbf{Reader's note on marking.} This comparative print shows the
changes H. R. 9510 makes to existing law, reproduced from the Federal Food, Drug, and
Cosmetic Act (21 U.S.C. 301 et seq.) as carried in Title 21, United States Code,
current through Public Law 119-93. Existing law in which no change is proposed is
shown in plain text. Matter proposed to be inserted is shown
\cprinsert{like this insert marker}; matter proposed to be deleted is shown
\cprstrike{like this strike marker}, with the marker text written out literally, so a
reader need not rely on strikethrough or underline rendering. Unchanged matter
omitted for brevity is shown by a row of asterisks. The print reaches the eleven
affected Title 21 device sections, reproduced in 21 U.S.C. order and shown only as to
their affected provisions, consistent with section 4 of the Act. It carries no new
duty; it is the transparency record a committee expects before introduction
\cite{pl-119-93, usc-fdca}.}

\uscprov{1}{\textbf{Section 301 (21 U.S.C. 331 note); Short Title.} A new statutory
note is added after the existing short-title notes of the Act:}
\uscquote{0}{Short Title of 2026 Amendment. \cprinsert{Pub. L. 119-XXX, section 1,
approved in 2026, provided that: ``This Act may be cited as the `Verification Before
Generation in Physical AI Oncology Trials Act of 2026'.''}}

\uscprov{1}{\textbf{Section 201(h) (21 U.S.C. 321(h)); Definition of Device.} The
device definition is amended by adding at the end of paragraph (1):}
\uscquote{0}{The term ``device'' ... does not include software functions excluded
pursuant to section 360j(o) of this title. \cprinsert{A software function that
generates or executes robot-patient interaction code (as defined in section 515D), or
that directs the autonomous physical action of a robot in contact with a human
subject in a clinical investigation, is a device and is not a software function
excluded under section 360j(o).}}

\uscprov{1}{\textbf{Section 301 (21 U.S.C. 331); Prohibited Acts.} A new prohibited
act is added after subsection (iii):}
\uscquote{0}{(iii) The refusal or failure to follow an order under section 364g of
this title. \cprinsert{(jjj) The generation or execution of robot-patient interaction
code in a Physical AI oncology clinical investigation in violation of section 515D, or
the failure to make or file the documentation and attestation required by section
515D before such generation or execution.}}

\uscprov{1}{\textbf{Section 501 (21 U.S.C. 351); Adulterated Devices.} A new
paragraph is added after subsection (j):}
\uscquote{0}{(j) Delayed, denied, or limited inspection ... refuses to permit entry
or inspection. \cprinsert{(k) Physical AI device software function used without a
cleared verification record. If it is a Physical AI device software function (as
defined in section 515D) used to generate or execute robot-patient interaction code in
a clinical investigation without a cleared verification record required by section
515D, or in a manner not in conformity with a predetermined change control plan
governed by section 515D.}}

\uscprov{1}{\textbf{Section 505F (21 U.S.C. 355g); Utilizing Real World Evidence.}
Subsection (b) is amended, and a new subsection (g) is added:}
\uscquote{0}{(b) Real world evidence defined. ... means data regarding the usage,
or the potential benefits or risks, of a drug derived from sources other than
traditional clinical trials. \cprinsert{Such term includes verification, validation,
and uncertainty quantification records and trial telemetry generated under section
515D.}}
\uscquote{0}{\cprinsert{(g) Application to Physical AI oncology investigations. The
Secretary shall, in carrying out the program under this section, evaluate the use of
real world evidence, including records generated under section 515D, to support the
evaluation of devices used in Physical AI oncology clinical investigations.}}

\uscprov{1}{\textbf{Section 510(k) (21 U.S.C. 360(k)); Premarket Notification.} The
subsection is amended by adding at the end:}
\uscquote{0}{... Such certification shall not be considered an element of such
notification. \cprinsert{A change to a device that generates or executes robot-patient
interaction code, cleared under this subsection, that is consistent with a
verification protocol under section 515D and an established predetermined change
control plan under section 360e-4, shall not require a new premarket notification under
this subsection.}}

\uscprov{1}{\textbf{Section 513 (21 U.S.C. 360c); Classification of Devices.} A new
subsection (l) is added:}
\uscquote{0}{\cprinsert{(l) Physical AI device software functions. In classifying a
device that generates or executes robot-patient interaction code, the Secretary shall
consider the degree of autonomy of the device (assistive, augmentative, or
autonomous) and may establish special controls, including a verification protocol
under section 515D, commensurate with that degree of autonomy. A device of autonomous
degree used in a Physical AI oncology clinical investigation shall be subject to the
verification requirements of section 515D regardless of class.}}

\uscprov{1}{\textbf{Section 515 (21 U.S.C. 360e); Premarket Approval.} The contents
of a premarket approval application are amended by striking the conjunction at the end
of subparagraph (G), revising subparagraph (H), and adding new subparagraph (I):}
\uscquote{0}{(G) the certification required under section 282(j)(5)(B) of title 42
(which shall not be considered an element of such application)\cprstrike{; and}}
\uscquote{0}{(H) such other information relevant to the subject matter of the
application as the Secretary, with the concurrence of the appropriate panel under
section 360c of this title, may require\cprstrike{.}\cprinsert{; and}}
\uscquote{0}{\cprinsert{(I) in the case of a device that generates or executes
robot-patient interaction code, the cleared verification record and the gate evidence
required by section 515D.}}

\uscprov{1}{\textbf{Section 515C (21 U.S.C. 360e-4); Predetermined Change Control
Plans.} The scope paragraphs for approved and cleared devices are amended, and a new
subsection (d) is added. In subsection (a)(3) (approved devices), and identically in
subsection (b)(3) (cleared devices), the scope clause is amended by inserting before
the period:}
\uscquote{0}{... and performance requirements for changes made under the
plan\cprinsert{, and, for a device that generates or executes robot-patient
interaction code, an automated verification-before-change protocol consistent with
section 515D that must clear before any change is implemented}.}
\uscquote{0}{\cprinsert{(d) Physical AI device software functions. A predetermined
change control plan for a device that generates or executes robot-patient interaction
code shall require that each planned change clear the verification process required by
section 515D, in advance of generation or execution, as a condition of being
implemented without a new submission under section 360e or section 360(k).}}

\uscprov{1}{\textbf{Section 520(o) (21 U.S.C. 360j(o)); Clinical Decision Support
Software.} Paragraph (1) is amended by adding at the end:}
\uscquote{0}{... so that it is not the intent that such health care professional rely
primarily on any of such recommendations to make a clinical diagnosis or treatment
decision regarding an individual patient. \cprinsert{The exclusion in this paragraph
does not apply to, and the term ``device'' includes, a software function that
generates or executes robot-patient interaction code, or that directs the autonomous
physical action of a robot in contact with a human subject in a clinical
investigation; such a function is a device subject to section 515D.}}

\uscprov{1}{\textbf{Section 521 (21 U.S.C. 360k); State and Local Requirements.} A
new subsection (c) is added:}
\uscquote{0}{\cprinsert{(c) Rule of construction for Physical AI oversight. Nothing
in this section preempts a State requirement of general applicability that a licensed
human retain the authority to monitor, intervene, override, or terminate an artificial
intelligence or robotic system, or that an audit record be maintained, to the extent
that the requirement is not different from, or in addition to, a requirement
applicable under this chapter to a device, including the verification requirements of
section 515D.}}

\uscprov{1}{\textbf{New Section 515D (21 U.S.C. 360e-5); Verification Before
Generation.} The following new section is inserted in subchapter V immediately after
section 515C (21 U.S.C. 360e-4). It is added in its entirety; the full text of
subsections (a) through (j) is set out in section 3 of this Act:}
\uscquote{0}{\cprinsert{Sec. 515D. Verification before generation of robot-patient
interaction code in Physical AI oncology investigations. (a) Requirement. No
robot-patient interaction code may be generated or executed in a Physical AI oncology
clinical investigation unless an automated verification, validation, and uncertainty
quantification process, bound to one or more named external standards, has first
cleared for that interaction, and the cleared verification record has been documented
and attested under subsection (e). (b) through (j) Order of operations, gate
thresholds and decision rule, readiness gates, documentation and attestation,
cybersecurity, human oversight, and lifecycle, nondiscrimination, regulations,
definitions, and rule of construction, are set out in full in section 3 of this Act.}}

\uscprov{1}{Each change above is an insertion except the two short strikes in section
515 (the conjunction at the end of subparagraph (G) and the period at the end of
subparagraph (H)), so existing law is otherwise preserved in full and the new rule is
threaded consistently through Title 21 \cite{usc-360e, usc-360e4}.}
